\chapter{Pop-11 fundamentals}

Pop-11 reads like an Algol-descended imperative language and behaves like a
Lisp with a stack machine underneath: infix syntax, block keywords that
always close, dynamic typing, garbage collection, first-class procedures, an
extensible reader. What makes it distinctive is the \emph{open stack}, and
almost everything idiomatic follows from it. The examples in this chapter
are the pieces of \ghfile{fleet.p}, the core of the Robot Army.

\section{Syntax at a glance}

\begin{pop}
;;; comments run from three semicolons to end of line
vars charge = 100;            ;;; a global (permanent) variable
charge - 10 -> charge;        ;;; ASSIGNMENT POINTS RIGHT: value -> destination
npr('charge ' sys_>< charge); ;;; print a line; sys_>< concatenates anything
charge =>                     ;;; the REPL's printer: "** 90"
\end{pop}

The assignment arrow is the first thing to internalise: \code{-> x} reads
``into \code{x}''. It points that way because assignment is a stack
operation --- the value is already there, and the arrow says where to put it.

\section{Words and strings}

\begin{pop}
'advance to ridge'   ;;; single quotes: a mutable vector of characters
"r1"                 ;;; double quotes: a WORD -- an interned symbol, not a string
\end{pop}

This is the single most common beginner error, and it has a silent cousin.
Words are interned: every \code{\dq r1\dq} is the same object. Strings are
not: two \code{'r1'}s are two objects that merely print alike. Poplog's
default hash table, \code{newproperty}, matches keys by \emph{identity}, so a
registry keyed by unit id must use words:

\begin{pop}
vars fleet = newproperty([], 64, false, true);
90 -> fleet("r1");
fleet("r1") =>            ;;; ** 90
fleet('r1') =>            ;;; ** <false>  -- a different object, silently
\end{pop}

Use words as keys, or use \code{newmapping}, which compares with \code{=}.
Both are right; silently missing is not.

\section{Lists, and the quoting trap}

An order in the Robot Army is a list of words, and list brackets
\emph{quote} their contents, like Lisp's \code{'(...)}:

\begin{pop}
vars order = [unit r1 advance to ridge];
order(2) =>               ;;; ** r1
\end{pop}

So \code{[unit u advance to target]} is a list of the literal words
\code{unit u advance to target}. To build an order from values, escape with
\code{\textasciicircum} (one item) or \code{\textasciicircum\textasciicircum}
(splice a list):

\begin{pop}
vars u = "r4", target = "ridge", squad = [r1 r2];
[unit ^u advance to ^target] =>     ;;; ** [unit r4 advance to ridge]
[all ^^squad hold] =>               ;;; ** [all r1 r2 hold]
\end{pop}

Vectors, \code{\{1 2 3\}}, are the constant-time-subscript counterpart.

\section{Procedures}

A procedure's result is whatever it leaves behind; there is no \code{return}
statement, and the last expression is the value. A procedure may leave
\emph{several} results, which the caller collects in reverse order:

\begin{pop}
define split_order(o); hd(o), tl(o) enddefine;
vars verb, rest;
split_order([advance to ridge]) -> rest -> verb;   ;;; verb = advance, rest = [to ridge]
\end{pop}

Closures are made by partial application with \code{\%\ldots\%}:

\begin{pop}
define drain(charge, n); charge - n enddefine;
vars step_cost = drain(% 2 %);      ;;; every step costs 2%
step_cost(100) =>                   ;;; ** 98
\end{pop}

And every accessor has an \emph{updater} --- the procedure that runs when it
appears on the left of an arrow. That is why \code{roster(2)} both reads and
writes:

\begin{pop}
vars roster = [r1 r2 r3];
"r9" -> roster(2);
roster =>                           ;;; ** [r1 r9 r3]
\end{pop}

\section{The open stack}

There is one user-visible stack, and the calling convention does not hide
it. A procedure can leave zero, one, or a hundred results, and the number
need not be known until it runs. \code{fleet.p} uses this for the roster:
\code{units()} walks the registry and leaves every robot on the stack, and
the \emph{caller} chooses the container.

\begin{pop}
define units();
    appproperty(fleet, procedure(k, v); v endprocedure);
enddefine;

define muster();  [% units() %]  enddefine;       ;;; a list of every unit

define fit_for_duty();
    lvars r;
    [% for r in muster() do
           if rb_charge(r) >= 25 then r endif       ;;; leave it, or don't
       endfor %]
enddefine;
\end{pop}

\code{[\%\ \ldots\ \%]} means ``run this, and make a list of everything it
left''. That one construct replaces most of what other languages need
generators, varargs, spread operators and comprehensions for. It also makes
Forth nearly free to host (\S5.4): Forth's data stack \emph{is} this stack.

\begin{callout}{Why this matters for performance}
Because results are passed on the stack rather than boxed into tuples,
multiple return values cost nothing, and neither does a procedure that
sometimes returns nothing. The \code{nfib} and \code{closures1M} benchmark
rows in Chapter~1 are measuring this calling convention.
\end{callout}

\section{Control structures}

Every opener has a matching closer.

\begin{pop}
if charge < 25 then recall() elseif charge < 60 then light_duty() else deploy() endif;
unless linked then hail() endunless;
for r in muster() do ... endfor;    for i from 1 to n do ... endfor;
while cond do ... endwhile;         repeat ... quitif(done); ... endrepeat;

define readiness(charge);
    switchon charge
        case < 25 then 'return to base'
        case < 60 then 'light duty only'
        else 'fit for duty'
    endswitchon
enddefine;
readiness(20) =>                    ;;; ** return to base
\end{pop}

\section{Scope, and \code{dlocal}}

\code{vars} declares a permanent (global) identifier; \code{lvars} declares a
lexical local inside a procedure --- prefer \code{lvars}. \code{constant} and
\code{lconstant} are their immutable counterparts.

\code{dlocal} has no close equivalent elsewhere: it gives a variable a new
value for the \emph{dynamic extent} of a call and restores the old value on
exit --- whether the procedure returns, jumps out, or is unwound by a mishap.
Here an order is relayed down three echelons, and the indent follows it:

\begin{pop}
vars echelon = '';
define relay(order, depth);
    dlocal echelon = echelon >< '> ';
    npr(echelon sys_>< order);
    if depth > 1 then relay(order, depth - 1) endif;
enddefine;
relay('advance', 3);
\end{pop}
\begin{repl}
> advance
> > advance
> > > advance
\end{repl}

After \code{relay} returns, \code{echelon} is empty again. The same mechanism
is how Poplog rebinds output streams and the compiler's own state safely ---
it is Pop-11's answer to \code{try\ldots finally}, and it is declarative.

\section{The list matcher}

Pop-11's pattern matcher works on lists and is built into the language with
the \code{matches} operator: \code{?x} binds one item, \code{??x} binds a
segment, \code{=} matches one item without binding. It is what turns an
order into an action in \code{fleet.p}:

\begin{pop}
vars id, place, n;            ;;; pattern variables must be permanent (vars)

define obey(order) -> reply;
    lvars r;
    'unintelligible' -> reply;
    if order matches [unit ?id advance to ?place] then
        fleet(id) -> r;
        if r then
            place -> rb_place(r);
            rb_charge(r) - 10 -> rb_charge(r);
            'unit ' sys_>< id sys_>< ' advancing to ' sys_>< place -> reply
        else 'no such unit: ' sys_>< id -> reply endif
    elseif order matches [unit ?id recharge ?n] then
        ...
    elseif order matches [report] then
        '' sys_>< length(fit_for_duty()) sys_>< ' of ' sys_><
        length(muster()) sys_>< ' units fit for duty' -> reply
    endif;
enddefine;
\end{pop}
\begin{repl}
[report]  ->  3 of 4 units fit for duty
[unit r1 advance to ridge]  ->  unit r1 advancing to ridge
[unit r9 advance to ridge]  ->  no such unit: r9
\end{repl}

\begin{callout}{A trap worth naming}
The matcher assigns through the \emph{permanent} identifier, so pattern
variables declared with \code{lvars} are silently not filled in --- you get
\code{<undef>}. Declare them \code{vars}.
\end{callout}

The same matcher, plus a dozen lines of ranking and reassembly, is the whole
of ELIZA; \code{examples/eliza.p} in the tree is a faithful implementation of
Weizenbaum's original.

\section{Objects: Objectclass}

Objectclass is Pop-11's CLOS-like object system: classes with slots, and
generic procedures with methods dispatched on argument class. A robot is one:

\begin{pop}
uses objectclass;

define :class Robot;
    slot rb_id     = "unknown";
    slot rb_kind   = "scout";      ;;; scout | sapper | medic | hauler
    slot rb_charge = 100;
    slot rb_place  = "base";
enddefine;

define :method print_instance(r:Robot);
    printf('<%p %p  %p%%  at %p>\n',
           [% rb_kind(r), rb_id(r), rb_charge(r), rb_place(r) %]);
enddefine;

define enlist(id, kind) -> r;
    newRobot() -> r;
    id -> rb_id(r);  kind -> rb_kind(r);
    r  -> fleet(id);
enddefine;
\end{pop}
\begin{repl}
<scout r1  90%  at base>
<sapper r2  100%  at base>
\end{repl}

Slots come with accessors \emph{and} updaters for free, which is why
\code{rb\_charge(r)} works on both sides of the arrow in \code{obey}.
